\lysh{22508_SOLUTION}{1}
\textbf{ΛΥΣΗ 2}

α) Το μέσο $K$ του ευθύγραμμου τμήματος $A\Gamma$ έχει συντεταγμένες
\[
x_K = \frac{x_A+x_\Gamma}{2} = \frac{1+3}{2} = 2 \text{ και } y_K = \frac{y_A+y_\Gamma}{2} = \frac{4+0}{2} = 2 .
\]
Ο συντελεστής διεύθυνσης της ευθείας $A\Gamma$ είναι ίσος με $\lambda_1 = \frac{4-0}{1-3} = -2$, κατά συνέπεια η μεσοκάθετος του τμήματος $A\Gamma$ έχει συντελεστή διεύθυνσης $\lambda_2 = \frac{1}{2}$.
Άρα η εξίσωση της μεσοκαθέτου του τμήματος $A\Gamma$ είναι $y - 2 = \frac{1}{2}(x - 2)$.

β) Το κέντρο του κύκλου διαμέτρου $A\Gamma$ είναι το σημείο $K(2,2)$ και η ακτίνα είναι ίση με $\rho = (AK) = \sqrt{(1 - 2)^2 + (4 - 2)^2} = \sqrt{5}$.
Επομένως η εξίσωση του κύκλου διαμέτρου $A\Gamma$ είναι $(x - 2)^2 + (y - 2)^2 = 5$.

γ) Οι ζητούμενες κορυφές $B,\Delta$ του τετραγώνου $AB\Gamma\Delta$ ισαπέχουν από τα σημεία $A,\Gamma$ και βλέπουν το $A\Gamma$ υπό ορθή γωνία, άρα είναι τα σημεία τομής της μεσοκαθέτου του τμήματος $A\Gamma$ και του κύκλου διαμέτρου $A\Gamma$.

Λύνουμε το σύστημα :
\[
y - 2 = \frac{1}{2}(x - 2) \quad (1)
\]
\[
(x - 2)^2 + (y - 2)^2 = 5 \quad (2)
\]
Η εξίσωση (2) με τη βοήθεια της (1) γράφεται
\begin{align*}
(x - 2)^2 + \frac{1}{4}(x - 2)^2 &= 5 \\
\Leftrightarrow 4(x - 2)^2 + (x - 2)^2 &= 20 \\
\Leftrightarrow 5(x - 2)^2 &= 20 \\
\Leftrightarrow (x - 2)^2 &= 4 \\
\Leftrightarrow x - 2 = 2 \text{ ή } x - 2 &= -2 \\
\Leftrightarrow x = 4 \text{ ή } x &= 0 .
\end{align*}
Για $x = 4$ βρίσκουμε $y = 3$ και για $x = 0$ βρίσκουμε $y = 1$.
Επομένως οι ζητούμενες κορυφές είναι τα σημεία $(4,3)$ και $(0,1)$.
\end{lysh}